\documentclass{standalone}
\usepackage{tikz}
\usetikzlibrary{calc}

\begin{document}

\begin{tikzpicture}[
  scale=1.5, 
  font=\small,
  node distance=2cm,
  dot/.style={circle, fill=black, inner sep=1.5pt},
  level 1/.style={sibling distance=5cm},
  level 2/.style={sibling distance=2.5cm},
  level 3/.style={sibling distance=1.2cm}
]

% Узел Природы
\node[dot, label=above:{Природа}] (nature) {}
    child {
        node[dot, label=left:{Игрок 1}] (P1_A) {}
        edge from parent node[left, pos=0.5] {$p$ (Тип A: Агрессор)}
    }
    child {
        node[dot, label=right:{Игрок 1}] (P1_B) {}
        edge from parent node[right, pos=0.5] {$1-p$ (Тип B: Мирный)}
    };

% Ветви для Типа A
\node[dot] (P2_A_H) at ($(P1_A) + (-1.2,-1)$ ) {};
\node[dot] (P2_A_D) at ($(P1_A) + (1.2,-1)$ ) {};

\draw (P1_A) -- (P2_A_H) node[midway, left] {Ястреб (H)};
\draw (P1_A) -- (P2_A_D) node[midway, right] {Голубь (D)};

% Ветви для Типа B
\node[dot] (P2_B_H) at ($(P1_B) + (-1.2,-1)$ ) {};
\node[dot] (P2_B_D) at ($(P1_B) + (1.2,-1)$ ) {};

\draw (P1_B) -- (P2_B_H) node[midway, left] {H};
\draw (P1_B) -- (P2_B_D) node[midway, right] {D};

% Выигрыши для Типа A
\draw (P2_A_H) -- ++(-0.6,-0.8) node[below] {$(v-c, v-c)$} node[midway, left] {H};
\draw (P2_A_H) -- ++(0.6,-0.8) node[below] {$(v, 0)$} node[midway, right] {D};
\draw (P2_A_D) -- ++(-0.6,-0.8) node[below] {$(0, v)$} node[midway, left] {H};
\draw (P2_A_D) -- ++(0.6,-0.8) node[below] {$(v/2, v/2)$} node[midway, right] {D};

% Выигрыши для Типа B (например, c=0 для мирного типа)
\draw (P2_B_H) -- ++(-0.6,-0.8) node[below] {$(v, v)$} node[midway, left] {H};
\draw (P2_B_H) -- ++(0.6,-0.8) node[below] {$(v, 0)$} node[midway, right] {D};
\draw (P2_B_D) -- ++(-0.6,-0.8) node[below] {$(0, v)$} node[midway, left] {H};
\draw (P2_B_D) -- ++(0.6,-0.8) node[below] {$(v/2, v/2)$} node[midway, right] {D};

% Информационное множество Игрока 2
\draw[dashed, rounded corners=10pt] 
    ($(P2_A_H) + (-0.2,0.2)$) rectangle ($(P2_B_H) + (0.2,-0.2)$);
\node at ($(P2_A_H)!0.5!(P2_B_H) + (0, 0.4)$) {Инф. множество Игрока 2};

\draw[dashed, rounded corners=10pt] 
    ($(P2_A_D) + (-0.2,0.2)$) rectangle ($(P2_B_D) + (0.2,-0.2)$);

\end{tikzpicture}
\end{document}
